\section{Method}
\label{sec:method}
GNN-QMIX has two networks (Fig.~\ref{fig:architecture}). The agent network is a GCN that
every vehicle runs on its own observation graph to obtain an embedding $h_i$ and its
individual action values $Q_i$. The mixer is a GAT over the graph of agents that is used
only during training: it reads the embeddings $h_i$ together with the global entity state
and generates the weights of a monotonic mixing function. The GAT makes the mixer
independent of the number of agents; the GCN enriches each agent's observation before it
reaches the mixer. Execution is decentralized: vehicle $i$ selects
$a_i = \arg\max_{a} Q_i(o_i,a)$ from its own observation graph, and no network ever selects
a joint action. Sections~\ref{subsec:entities}--\ref{subsec:training} specify stage~1 of
the work (GNN-QMIX, learned end to end). Section~\ref{subsec:hybrid} describes stage~2
(PID + GNN-QMIX).

\begin{figure}[hbt!]
\centering
\fbox{\parbox[c][0.18\textheight][c]{0.9\textwidth}{\centering
Figure placeholder: (a) GCN agent on the sensor-radius observation graph;
(b) overall CTDE architecture; (c) monotonic GAT mixer.}}
% TODO(authors): draw the architecture figure and replace the box with
% \includegraphics[width=0.9\textwidth]{architecture}
\caption{\label{fig:architecture} Architecture of GNN-QMIX. Each vehicle encodes its
observation graph with a shared GCN to obtain $h_i$ and $Q_i$. During training, a GAT over
the agent graph takes $[\mathbf{S}_i \,\Vert\, h_i]$ for every agent and generates the
non-negative mixing weights $\mathbf{W}_1$, $\mathbf{W}_2$ and the biases $\mathbf{b}_1$,
$b_2$ that combine the chosen $Q_i$ into $Q_{\text{tot}}$.}
\end{figure}

\subsection{Graph Observations and State}
\label{subsec:entities}
Following TransfQMix~\cite{transfqmix}, observations and the state are matrices with one row
per entity and one column per feature type, instead of concatenated vectors. Features of
the same type are then processed by the same weights for every entity, and the number of
rows may change without changing any weight shape. The only entities in the present
problem are the cooperating vehicles, and all features are kinematic.

\textbf{Observation graph.} Vehicle $i$ senses the vehicles inside its sensor radius,
\begin{equation}
\label{eq:neighbors}
\mathcal{N}_i = \{\, j \neq i \;:\; d_{ij} \le r_{\text{detect}} \,\}.
\end{equation}
Its observation graph $G_i=(V_i,E_i)$ has the node set $V_i=\{i\}\cup\mathcal{N}_i$. A vehicle
is added to $G_i$ when it comes inside the sensor radius of vehicle $i$ and is removed when
it leaves. Vehicle $i$ is joined to every node in $\mathcal{N}_i$. Two neighbors
$j,k\in\mathcal{N}_i$ are joined when $d_{jk}\le r_{\text{detect}}$, which vehicle $i$ can
compute from the two relative positions it measures. Every node has a self-loop. In contrast
to TransfQMix, where the agent attends over all observable entities and learns the edge
weights, the edges here are fixed by the sensor geometry.

\textbf{Node features.} Each node carries 11 features, arranged as
$[\,\text{own}\ (3)\;|\;\text{relative}\ (7)\;|\;\texttt{IS\_SELF}\ (1)\,]$. The self node
holds the own features of vehicle $i$,
\begin{equation}
\label{eq:own}
\big(\cos\chi_i,\ \sin\chi_i,\ v_i\big),
\end{equation}
where $\chi_i$ is the drift angle, that is, the difference between the heading of vehicle
$i$ and the bearing to its waypoint, and $v_i$ is its normalized airspeed; it has
\texttt{IS\_SELF}~$=1$ and zeros in the relative block. The node of a neighbor $j$ holds the
relative features measured by vehicle $i$,
\begin{equation}
\label{eq:rel}
\big(\Delta x_{ij},\ \Delta y_{ij},\ \Delta v^{x}_{ij},\ \Delta v^{y}_{ij},\
\cos\psi_{ij},\ \sin\psi_{ij},\ d_{ij}\big),
\qquad \psi_{ij} = \operatorname{atan2}\!\big(\Delta v^{y}_{ij},\,\Delta v^{x}_{ij}\big),
\end{equation}
with positions and velocities expressed in a local east--north frame and divided by fixed
scales; it has \texttt{IS\_SELF}~$=0$ and zeros in the own block. The \texttt{IS\_SELF} flag
is the one introduced by TransfQMix to mark the observing agent after the positional
information of a concatenated vector is lost. The observation of vehicle $i$ is therefore
the pair of an entity matrix and an adjacency matrix,
\begin{equation}
\label{eq:obs-matrix}
\mathbf{O}_i \in \mathbb{R}^{(1+|\mathcal{N}_i|)\times 11},
\qquad
\mathbf{A}_i \in \{0,1\}^{(1+|\mathcal{N}_i|)\times(1+|\mathcal{N}_i|)} .
\end{equation}

\textbf{State.} The global state used by the mixer is the entity matrix
\begin{equation}
\label{eq:state-matrix}
\mathbf{S} \in \mathbb{R}^{n\times 7},
\qquad
\mathbf{S}_i = \big(x_i,\ y_i,\ v^{x}_i,\ v^{y}_i,\ \cos\chi_i,\ \sin\chi_i,\ v_i\big),
\end{equation}
where $(x_i,y_i)$ is the position of vehicle $i$ relative to the centroid of the fleet and
$(v^x_i,v^y_i)$ is its ground velocity, with the same scaling as in the observation. The
state has no \texttt{IS\_SELF} feature because it is not defined with respect to an agent.

\subsection{GCN Agent}
\label{subsec:agent}
The agent network is shared by all vehicles. It embeds the node features with a linear
layer, applies $L_{\text{GCN}}$ graph convolutional layers of the form of Eq.~\eqref{eq:gcn}
on the observation graph, and reads the embedding of the self node:
\begin{align}
\label{eq:agent-embed}
\mathbf{H}_i^{(0)} &= \mathrm{ReLU}\big(\mathbf{O}_i\,\mathbf{W}_{\text{emb}}\big),
\qquad \mathbf{W}_{\text{emb}}\in\mathbb{R}^{11\times d_h},\\
\label{eq:agent-gcn}
\mathbf{H}_i^{(\ell+1)} &= \mathrm{ReLU}\Big(\mathbf{D}_i^{-1/2}\,\mathbf{A}_i\,\mathbf{D}_i^{-1/2}\,
\mathbf{H}_i^{(\ell)}\,\mathbf{W}^{(\ell)}\Big),
\qquad \ell = 0,\dots,L_{\text{GCN}}-1,\\
\label{eq:agent-q}
h_i &= \big[\mathbf{H}_i^{(L_{\text{GCN}})}\big]_{\text{self}},
\qquad
Q_i(o_i,\cdot) = \mathrm{MLP}_{Q}(h_i) \in \mathbb{R}^{|\mathcal{A}|}.
\end{align}
Here $\mathbf{A}_i$ already contains the self-loops and $\mathbf{D}_i$ is its degree matrix;
bias terms are omitted. The embedding $h_i\in\mathbb{R}^{d_h}$ summarizes the local traffic
situation of vehicle $i$ after $L_{\text{GCN}}$ rounds of message passing among the
vehicles it senses. A vehicle outside the sensor radius is not a node of $G_i$ and cannot
influence $Q_i$. The agent has no recurrent input: unlike the TransfQMix agent, it carries
no hidden state between time steps.
% TODO(authors): deferred per decision 2026-11-25 (GAP-07/GAP-09/GAP-12) -- no GRU hidden
% state enters the node features; recurrence belongs to the latency work.

\subsection{Monotonic GAT Mixer}
\label{subsec:mixer}
The mixer operates on a graph with one node per agent. The input feature of node $i$ is the
concatenation of its state row and its agent embedding,
\begin{equation}
\label{eq:mixer-input}
x_i = \big[\,\mathbf{S}_i \,\Vert\, h_i\,\big] \in \mathbb{R}^{7+d_h},
\end{equation}
and gradients of the loss flow through $h_i$ into the agent network. In the default
configuration the agent graph is fully connected and the attention coefficients of
Eq.~\eqref{eq:gat-alpha} determine how much each agent contributes to another, as in the
TransfQMix mixer; an alternative configuration joins two agents only when
$d_{ij}\le r_{\text{detect}}$ on the true positions. The node features are embedded and
passed through $L_{\text{GAT}}$ multi-head graph attention layers
(Eqs.~\eqref{eq:gat-alpha}--\eqref{eq:gat}),
\begin{equation}
\label{eq:mixer-gat}
g_i^{(0)} = \mathrm{ReLU}\big(\mathbf{W}_{x}\,x_i\big),
\qquad
g_i^{(\ell+1)} = g_i^{(\ell)} + \mathrm{ELU}\Big(\big[\mathrm{GAT}^{(\ell)}\big(g^{(\ell)}_1,\dots,g^{(\ell)}_n\big)\big]_i\Big),
\qquad
g_i = g_i^{(L_{\text{GAT}})} \in \mathbb{R}^{d_h}.
\end{equation}
% NOTE: the residual connection and the ELU in Eq. (mixer-gat) follow
% algorithms/gnn_qmix/networks.py as read on 2026-11-25; the plan (G05.3 section 3.3) only
% states "gat_layers layers of multi-head GAT". TODO(authors): confirm.
The node embedding $g_i$ combines the local reasoning of agent $i$ with information the
agent did not have, namely the embeddings of the other agents and the true state. The
mixing terms of Eq.~\eqref{eq:qmix} are generated from the node embeddings as follows:
\begin{align}
\label{eq:w1}
\mathbf{W}_{1,i} &= \big|\,\mathrm{Linear}_{W_1}(g_i)\,\big| \in \mathbb{R}^{d_m}
&&\text{(one row per agent)},\\
\label{eq:readout}
p &= \Big[\;\frac{1}{n}\sum_{i=1}^{n} g_i \;\Big\Vert\; \max_{i} g_i\;\Big] \in \mathbb{R}^{2d_h}
&&\text{(fixed size for any } n\text{)},\\
\label{eq:b1w2b2}
\mathbf{b}_1 &= \mathrm{Linear}_{b_1}(p), \qquad
\mathbf{W}_2 = \big|\,\mathrm{Linear}_{W_2}(p)\,\big|, \qquad
b_2 = \mathrm{MLP}_{b_2}(p),
\end{align}
where the maximum in Eq.~\eqref{eq:readout} is taken elementwise over the agents and
$|\cdot|$ is the elementwise absolute value. The joint action value is
\begin{equation}
\label{eq:qtot}
Q_{\text{tot}} = \mathrm{ELU}\Big(\frac{1}{n}\sum_{i=1}^{n} Q_i(o_i,a_i)\,\mathbf{W}_{1,i} + \mathbf{b}_1\Big)\cdot\mathbf{W}_2 + b_2 .
\end{equation}
Equation~\eqref{eq:qtot} is the QMIX mixer of Eq.~\eqref{eq:qmix} with two changes. First,
row $i$ of $\mathbf{W}_1$ is produced from the node embedding of agent $i$ by a layer that
is shared by all agents, so $\mathbf{W}_1$ has $n$ rows without any layer of size $n$; this
is the mechanism of TransfQMix. Second, the sum over agents is replaced by a mean. The
terms $\mathbf{b}_1$, $\mathbf{W}_2$, and $b_2$ come from the readout $p$ instead of the
recurrent tokens that TransfQMix adds to its mixer graph. The readout uses the mean and the
maximum, not the sum, so its scale does not grow with $n$. With mean aggregation in
Eq.~\eqref{eq:qtot} and the team reward of Eq.~\eqref{eq:team-reward}, $Q_{\text{tot}}$
stays on one scale for every fleet size. Aggregation by sum is available as a configuration
option.

\subsection{Properties}
\label{subsec:properties}
The two properties below follow directly from Eqs.~\eqref{eq:agent-embed}--\eqref{eq:qtot}.
They restate theoretical claim~1 of Section~\ref{sec:intro}.

\textbf{Proposition 1 (monotonicity and IGM).}
For every fleet size $n\ge 1$, every state, and every joint observation, the mixer of
Eq.~\eqref{eq:qtot} satisfies $\partial Q_{\text{tot}}/\partial Q_i \ge 0$ for all $i$.
Consequently the joint action formed by the individual greedy actions
$a_i = \arg\max_{a}Q_i(o_i,a)$ maximizes $Q_{\text{tot}}$, that is, Eq.~\eqref{eq:igm} holds.

\textit{Proof sketch.} Let $\mathbf{z} = \frac{1}{n}\sum_i Q_i\,\mathbf{W}_{1,i} + \mathbf{b}_1$
and write $[\cdot]_k$ for the $k$-th component. The quantities $\mathbf{W}_{1,i}$,
$\mathbf{b}_1$, $\mathbf{W}_2$, and $b_2$ are functions of $\mathbf{S}$ and of
$h_1,\dots,h_n$, and each $h_i$ is a function of the observation $o_i$ only, so none of them
depends on the selected actions. For a fixed state and joint observation, the chain rule
gives
\begin{equation}
\label{eq:dqtot}
\frac{\partial Q_{\text{tot}}}{\partial Q_i}
= \frac{1}{n}\sum_{k=1}^{d_m} \mathrm{ELU}'\big([\mathbf{z}]_k\big)\,[\mathbf{W}_{1,i}]_k\,[\mathbf{W}_2]_k .
\end{equation}
The derivative of the ELU is positive, and $[\mathbf{W}_{1,i}]_k\ge 0$ and
$[\mathbf{W}_2]_k\ge 0$ because of the absolute values in Eqs.~\eqref{eq:w1}
and~\eqref{eq:b1w2b2}, so every term of the sum is non-negative. This holds for any $n$
because $n$ enters only through the positive factor $1/n$ and the number of rows
$\mathbf{W}_{1,i}$. A function that is non-decreasing in each argument $Q_i$ attains its
maximum over joint actions when each $Q_i(o_i,a_i)$ is maximized separately, which is the
argument used for QMIX~\cite{qmix}. The same derivation holds with aggregation by sum.

\textbf{Proposition 2 (fleet-size independence and permutation symmetry).}
(a) No parameter of the agent network or of the mixer has a dimension that depends on the
fleet size $n$ or on the number of neighbors $|\mathcal{N}_i|$, so the same parameters
$\theta$ define $Q_i$ and $Q_{\text{tot}}$ for every $n$. (b) $Q_i$ does not depend on the
order in which the neighbors of vehicle $i$ are listed. (c) $Q_{\text{tot}}$ does not depend
on the order in which the agents are listed.

\textit{Proof sketch.} (a) The learned parameters are $\mathbf{W}_{\text{emb}}$
($11\times d_h$), the GCN weights $\mathbf{W}^{(\ell)}$ ($d_h\times d_h$), the output
network $\mathrm{MLP}_Q$ (input $d_h$, output $|\mathcal{A}|$), the mixer embedding
$\mathbf{W}_x$ (input $7+d_h$, output $d_h$), the GAT parameters (functions of $d_h$ and
the number of heads), and the layers in Eqs.~\eqref{eq:w1} and~\eqref{eq:b1w2b2} (inputs
$d_h$ or $2d_h$, outputs $d_m$ or $1$). Each is applied either to one node at a time or to
the fixed-size readout $p$, and the operations that combine nodes (the normalized adjacency
product, the attention-weighted sum, the mean, and the maximum) are defined for any number
of nodes. (b) By the permutation equivariance of the GCN layer
(Section~\ref{subsec:gnn-background}), reordering the neighbor nodes reorders the rows of
$\mathbf{H}_i^{(\ell)}$ in the same way and leaves the content of the self-node row, which is
the only row that is read out, unchanged. (c) Reordering the agents by
$\mathbf{P}$ reorders the inputs $x_i$ and, by the equivariance of the GAT layers, the
embeddings $g_i$ and the rows $\mathbf{W}_{1,i}$ in the same way. The mean and the maximum
in Eq.~\eqref{eq:readout} and the mean in Eq.~\eqref{eq:qtot} are symmetric functions, so
$p$, $\mathbf{b}_1$, $\mathbf{W}_2$, $b_2$, and $Q_{\text{tot}}$ are unchanged.

Proposition~2 states that the model can be evaluated at any $n$; it does not state that the
values it produces at an unseen $n$ are accurate. Whether the learned policy remains
effective at larger fleet sizes is an empirical question, addressed in
Section~\ref{subsec:results-zeroshot}.

\subsection{Training Over Several Fleet Sizes}
\label{subsec:training}
Because the parameters are independent of $n$ (Proposition~2), one model is trained on
several fleet sizes at once. For each training size $n\in\mathcal{N}_{\text{train}}$ there is
one environment process and one replay buffer $\mathcal{D}_n$ of single transitions
$(\mathbf{S},\mathbf{o},\mathbf{a},r,\mathbf{S}',\mathbf{o}',e)$. In every iteration of the
training loop each environment advances by one step with $\epsilon$-greedy decentralized
action selection. In every update one batch is drawn from each buffer, the loss of
Eq.~\eqref{eq:td} is evaluated separately for each fleet size, and the losses are summed
before a single optimizer step:
\begin{equation}
\label{eq:joint-loss}
\mathcal{L}(\theta) = \sum_{n\in\mathcal{N}_{\text{train}}} \mathcal{L}_n(\theta),
\qquad
\mathcal{L}_n(\theta) = \mathbb{E}_{\mathcal{D}_n}\Big[\big(y - Q_{\text{tot}}(\mathbf{S},\mathbf{o},\mathbf{a};\theta)\big)^2\Big].
\end{equation}
Each batch contains transitions of a single fleet size, so no padding of the agent
dimension is needed. The target is the one-step target of Eq.~\eqref{eq:td},
\begin{equation}
\label{eq:target}
y = r + \gamma\,(1-e)\,Q_{\text{tot}}\big(\mathbf{S}',\mathbf{o}',\mathbf{a}^{\star};\theta^-\big),
\qquad
a^{\star}_i = \arg\max_{a} Q_i(o'_i,a;\theta^-),
\end{equation}
where the maximization over joint actions is carried out per agent, which is exact by
Proposition~1. The target parameters follow the online parameters by Polyak averaging,
$\theta^- \leftarrow \tau\,\theta + (1-\tau)\,\theta^-$. Bootstrapping stops only on
termination ($e=1$) and not when an episode is truncated by a time limit.

TransfQMix reports zero-shot transfer from models trained at one team size and a curriculum
over team sizes; the scheme above instead mixes the fleet sizes within every update.

\subsection{Hybrid Control: PID + GNN-QMIX}
\label{subsec:hybrid}
Stage~2 places the learned policy inside a hybrid controller. A nominal controller tracks
the waypoint of each vehicle, and the learned conflict-resolution policy takes over a
vehicle only while that vehicle has traffic inside its sensor radius. Both act through the
same interface: a heading command and a speed command to the simulator (\texttt{HDG} and
\texttt{SPD} in BlueSky).

\textbf{Nominal controller.} Each vehicle runs a PID controller $\pi_{\text{PID}}$ that
tracks its waypoint. At every control step the controller proposes a command
$a_i^{\text{PID}}$ consisting of a heading that turns the vehicle toward its waypoint and a
speed.
% TODO(authors): the PID law is not yet specified (next plan, "PID + GNN-QMIX", goals
% G05.5--G05.9). To decide: the error signal (the drift angle \chi_i of Eq. (eq:own) is the
% natural candidate); the gains K_P, K_I, K_D; the speed law (hold the nominal speed?);
% and whether the PID output is a continuous HDG/SPD command or is quantized to the same 15
% discrete actions as the learned policy (GAP-12).

\textbf{Activation rule.} The learned policy is active for vehicle $i$ when the observation
graph of that vehicle contains at least one neighbor,
\begin{equation}
\label{eq:activation}
c_i = \mathbf{1}\big[\,|\mathcal{N}_i| \ge 1\,\big],
\end{equation}
with $\mathcal{N}_i$ from Eq.~\eqref{eq:neighbors}. The command applied to the vehicle is
\begin{equation}
\label{eq:switching}
a^{\text{final}}_i =
\begin{cases}
\arg\max_{a} Q_i(o_i,a), & c_i = 1,\\[2pt]
a^{\text{PID}}_i, & c_i = 0.
\end{cases}
\end{equation}
The rule that adds a vehicle to an observation graph is therefore also the trigger of the
learned policy: the policy acts exactly when its input graph is non-trivial. The rule is
evaluated per vehicle from that vehicle's own sensor data, so the hybrid controller remains
decentralized and needs no communication. When the last neighbor leaves the sensor radius,
$c_i$ returns to zero and control returns to the PID controller.
% NOTE (confirmed by the authors 2026-11-25, D-031): the extended abstract switched on a learned safety score
% s^i < \theta with \theta calibrated from a control barrier function. By the 2026-11-25
% decision the switching rule is the sensor radius; the safety score is mentioned only as a
% possible extension in the conclusion.
% TODO(authors): the extended abstract supplied the vector of PID proposals a_PID to the
% learned policy as an input (so that it learns corrections relative to the PID). Whether the
% stage-2 policy receives a_i^PID (for example as an extra node feature) is not decided; the
% text above does not assume it.

\textbf{Intervention count.} The number of interventions in an episode of $T$ control steps
is the number of vehicle-steps in which the learned policy is active,
\begin{equation}
\label{eq:interventions}
I = \sum_{t=1}^{T}\sum_{i=1}^{n} c_i^{\,t},
\qquad 0 \le I \le nT ,
\end{equation}
and $I/(nT)$ is the intervention frequency.
% TODO(authors): confirm this definition of I (vehicle-steps with c_i = 1). The plan (G05.3
% section 3.5) lists intervention frequency as "not defined here: needs the PID layer".
The learner produces actions only for these $I$ vehicle-steps, while the remaining
$nT - I$ are handled by the fixed controller. Theoretical claim~2 of
Section~\ref{sec:intro} is the hypothesis that this restriction makes the learning problem
easier than in stage~1, where the policy must also learn waypoint tracking.

\textbf{Training in the hybrid setting.} The networks, the mixer, and the loss are those of
stage~1. Several details of how the hybrid controller interacts with the TD update are not
yet decided and are not assumed here.
% TODO(authors): decide how transitions with no active vehicle (c_i = 0 for all i) are used
% in training: discarded, or stored with the PID actions (GAP-15).
% TODO(authors): decide how Q_tot is formed when only some vehicles are active: mix over all
% vehicles with the PID action treated as the taken action, or over the active subset only.
% TODO(authors): decide whether an intervention cost is added to the reward (the extended
% abstract stated a cost per intervention) and its value.
% TODO(authors): the extended abstract planned to characterize the continuity of the handoff
% between learned commands and PID commands; decide whether this stays in scope.
